\labsection{5}{Design graphics that answer a question}{sec:lab05}

\begin{labproject}{Lab 5 visual communication}
\textbf{Datasets.} \texttt{red\_wine\_quality\_dataset.csv}, with a Bike Sharing hourly extension.\\
\textbf{Goal.} Match chart form and visual encoding to the analytical task.

\textbf{Procedure.}
\begin{enumerate}
  \item Use a bar chart to show the frequency of each discrete quality score.
  \item Use side-by-side boxplots to compare alcohol content across quality levels.
  \item Use a scatterplot to inspect the joint behaviour of volatile acidity and alcohol.
  \item Redesign one chart by changing the question rather than merely changing colours or fonts.
  \item For each figure, write a caption that states the data, the encoding, the main visible pattern, and one limitation.
  \item Identify one misleading visual choice that would make the same data harder to interpret.
\end{enumerate}

\textbf{Interpretation checkpoint.} Quality score is ordered but discrete. Treating it as continuous may be convenient for some models, but a graphic should not hide that measurement scale.

\textbf{Extension visual.} Use \texttt{bike\_sharing\_dataset\_hour.csv} to create an hour-by-weekday heatmap of mean rentals. Explain what the heatmap reveals that a single hourly line chart does not, and what information is lost through aggregation.

\lstinputlisting[style=dsr]{all_experiments/lab05_visualization.R}
\end{labproject}
